% titel: Blatt 0, A10: Aussagen und ihre Gegenteile

% Teil I: Werkzeuge, kein Lösungsweg.

% Grundidee: Negation = "Umdrehen" der Quantoren und des Prädikats. Ein
% Allquantor ("alle") wird zur Existenzaussage, eine Existenzaussage
% ("es gibt") zu "keine". Nicht über Wahrheit entscheiden.

Die Negation einer Aussage ist die Aussage, die \emph{genau dann wahr}
ist, wenn die ursprüngliche falsch ist. Ihre Konstruktion folgt festen
Mustern, die nichts mit dem Wahrheitsgehalt zu tun haben --- es geht
darum, die Aussage ``\emph{umzudrehen}``.

\begin{werkzeugkasten}[Negation von Quantoren]
Die beiden wichtigsten Regeln:
\begin{align}
  \neg\bigl(\forall x: P(x)\bigr) &\Longleftrightarrow \exists x:
  \neg P(x),\\
  \neg\bigl(\exists x:P(x)\bigr) &\Longleftrightarrow \forall x:\neg P(x).
\end{align}
\textbf{Warum ``alle'' zu ``mindestens eine Gegenbeispiel`` wird:} Um eine
``alle''-Aussage zu widerlegen, genügt \emph{ein} Gegenbeispiel. Das
Gegenteil von ``jedes Tier ist grau`` ist also ``\emph{es gibt} ein Tier,
das nicht grau ist`` --- nicht ``kein Tier ist grau`` (viel zu stark).
Umgekehrt: ``es gibt eine`` wird zu ``\emph{keine}``.

\subsection*{Also: Allquantor $\leftrightarrow$ Existenzquantor, Prädikat
verneint.}
\end{werkzeugkasten}

\begin{werkzeugkasten}[Negation von Verknüpfungen (De Morgan)]
Logische ``und``/``oder`` werden beim Negieren ebenfalls umgedreht:
\begin{equation}
  \neg(P \text{ und } Q) \Longleftrightarrow (\neg P)\ \text{oder}\ (\neg Q).
\end{equation}
Beispiel: ``Es gibt eine Pflanze, die weder grün noch gelb ist`` ist
$\exists x:\ \neg(\text{grün }\text{oder}\ \text{gelb})$; die Negation
ist $\forall x:\ \neg\neg(\text{grün oder gelb}) =
\text{`jede Pflanze ist grün oder gelb'}$.
\end{werkzeugkasten}

\begin{achtung}[Typische Fehler]
\begin{itemize}
  \item \textbf{Negation einer ``alle``-Aussage} ist ``\emph{ein}
        Gegenbeispiel``, nicht ``keine``. ``Kein Tier ist grau`` ist nicht
        das Gegenteil von ``jedes Tier ist grau``.
  \item \textbf{Negation einer ``es gibt``-Aussage} ist ``keine``, nicht
        ``alle ausser...``.
  \item `P und Q` negiert man zu `$\neg P$ oder $\neg Q$` (nicht zu
        `$\neg P$ und $\neg Q$``).
  \item Nicht über Wahrheit entscheiden --- nur die Negation bilden.
\end{itemize}
\end{achtung}
